\title{Byte Latent Transformer: Patches Scale Better Than Tokens}

\begin{abstract}
We present the Byte Latent Transformer (\textsc{BLT}), an innovative LLM architecture operating directly at the byte level that, for the first time, achieves parity with tokenization-based LLMs in performance at scale, while offering notable gains in both inference robustness and computational efficiency. \textsc{BLT} transforms byte sequences into variably sized patches, utilizing these patches as the main computational elements. The model partitions these patches according to the estimated entropy of subsequent bytes, thereby distributing greater computational resources and model capacity to segments exhibiting higher data complexity. We perform the inaugural scaling analysis of byte-level language models with controlled floating point operations (flops), extending up to 8B model parameters and 4T bytes of training data. Our empirical findings validate that large-scale models can be efficiently trained directly from raw byte data, obviating the need for pre-defined token vocabularies. Both learning and inference procedures benefit from the adaptive selection of extended patches in regions of predictable data, yielding both computational improvements and enhanced performance on reasoning as well as rare cases requiring strong generalization. In summary, under equivalent inference budgets, \textsc{BLT} achieves more effective scaling compared to traditional tokenized models by enabling concurrent growth in both patch length and model size.
\end{abstract}

\section{Introduction}
\label{section:intro}

We present the Byte Latent Transformer~(\textbf{BLT}), an architecture that dispenses with tokenization by learning directly from raw bytes. Remarkably, BLT is the first model of its kind to achieve performance on par with traditional tokenization-based approaches at scale, while offering notable gains in both computational efficiency and robustness (\S\ref{sec:robustness}). Current large language models (llms) utilize an almost entirely end-to-end training pipeline, but maintain tokenization as a heuristic preliminary phase, segmenting bytes into a predetermined, fixed vocabulary. This reliance on static tokens introduces inherent biases in data compression, which in turn gives rise to known issues, including modality and domain sensitivity~\citep{dagangetting}, increased vulnerability to noisy or corrupted inputs~(\S\ref{sec:robustness}), insufficient capture of orthographic detail~\citep{edman2024cute}, and persistent challenges in supporting linguistic equity across multiple languages~\citep{liang2023xlm,petrov2024language,limisiewicz2024myte}.

Historically, tokenization has been a crucial component because training llms directly on byte-level data has been infeasible at large scales, primarily due to the excessive sequence lengths involved~\citep{xue2022byt5}. To address this, earlier research has turned to more efficient forms of self-attention~\citep{el2020characterbert, clark2022canine} or architectures that eliminate attention altogether~\citep{wang2024mambabyte}~(\S\ref{section:related_work}). Nonetheless, such approaches mainly benefit the training of \textit{small models}. When scaling up, the main computational bottleneck of Transformers becomes the expansive feed-forward network layers, which are executed for each individual byte, overshadowing the resource demands of the attention layers.

We introduce a dynamic and trainable approach for aggregating bytes into \textit{patches}~(\S\ref{section:patching}), along with a novel model architecture that integrates both byte-level and patch-level information. In contrast to tokenization, {\textsc BLT} does not rely on a predetermined vocabulary for patches. Instead, it utilizes efficient, learnable encoder and decoder components to translate any sequence of bytes into latent patch representations. Our results demonstrate that this method enables a \textit{greater} efficiency in compute allocation compared to models that depend on tokenization.

Tokenization-based language models distribute computational resources evenly across all tokens, regardless of each token's complexity. This approach prioritizes performance at the expense of efficiency, as the compression rules used for token induction do not consistently align with the intricacies of the prediction task. The foundation of our approach is the principle that computational effort should be allocated adaptively based on the demands of specific input regions. For instance, deploying a large transformer is unnecessary for most word suffixes, since such predictions usually involve low uncertainty and are straightforward, unlike the often ambiguous task of selecting the initial word in a sentence. This philosophy is embodied in the architecture of {\textsc BLT}~(\S\ref{section:architecture}), which integrates three transformer modules: two compact, byte-level \textit{local models} and one extensive global \textit{latent transformer}~(\autoref{fig:arch_high_level}). 

To inform decisions regarding the aggregation of bytes into patches—and thereby enable adaptive compute allocation—{\textsc BLT} segments sequences according to the entropy of the next-byte prediction, resulting in byte groupings that exhibit consistent information density across local contexts.

We introduce the inaugural flop-controlled scaling analysis of byte-level models, extending up to 8B parameters and 4T training bytes. Our results demonstrate that it is feasible to train large-scale models directly from bytes end-to-end, bypassing the need for fixed-vocabulary tokenization. In summary, {\textsc BLT} achieves training flop-controlled performance\footnote{The computational expenditure for each model is measured in terms of Floating Point OPerations (flops).} on par with Llama 3, yet requires as much as 50\% fewer inference flops~(\S\ref{section:scaling}).

Moreover, we find that processing raw bytes directly brings marked advancements in modeling data from the distribution’s long tail. {\textsc BLT} models outperform those relying on tokenization with greater resilience to input noise and show superior character-level capabilities, as evidenced in tasks such as orthographic analysis, phonological processing, and machine translation for low-resource languages~(\S\ref{sec:robustness}).

Additionally, {\textsc BLT} enables the simultaneous increase of both model and patch sizes without exceeding a fixed inference flop threshold. Utilizing larger patch sizes typically reduces computational overhead, freeing up resources to expand the global latent transformer since it needs to be activated less frequently. Through inference flop-controlled scaling studies~(\autoref{fig:fixed_inference_scaling}), we observe substantially improved scaling behaviors relative to those found in tokenization-based model architectures.

To summarize, this work presents several key advancements: 1) The proposal of {\textsc BLT}, a byte latent llm framework that adaptively assigns computational resources for enhanced flop efficiency; 2) Evidence that our approach achieves parity with Llama 3 in terms of flop-controlled training up to an 8B parameter scale, with the ability to exchange small drops in evaluation scores for up to 50\% improvements in flop efficiency; 3) The {\textsc BLT} architecture introduces a novel approach to scaling llms, enabling increases in model size while keeping inference costs constant; 4) We provide empirical results showing that {\textsc BLT} models exhibit greater resilience to noisy input and demonstrate sensitivity to sub-word information that token-based llms tend to overlook. The code for training and inference with {\textsc BLT} is publicly available at~\url{https://github.com/facebookresearch/blt}.

\section{Patching: From Individual Bytes to Groups of Bytes}
\label{section:patching}

Dividing bytes into discrete \textit{patches} enables {\textsc BLT} to flexibly assign computational resources depending on the input context. Figure~\ref{fig:patching_types} illustrates multiple approaches to grouping bytes into patches. More precisely, a patching function $f_p$ partitions a byte sequence $\pmb{x}=\{x_i,|i=1,\ldots n\}$ of length $n$ into $m < n$ patches, $\pmb{p}=\{p_j|j=1,\ldots,m\}$, by assigning each $x_i$ a label in \{0,1\}; a value of 1 marks the beginning of a new patch. In both token-based and patch-based frameworks, the main determinant of computation lies in the count of Transformer steps required, which in {\textsc BLT} corresponds to the number of patches generated via a particular patching function. Thus, the typical patch length (or simply \textit{patch size}) becomes the principal variable affecting computational demands—during both model training and evaluation—when a specific patching strategy is used~(\S\ref{section:flops}).  
In the following sections, we introduce three patching strategies: using a fixed number of bytes per patch~(\S\ref{section:static-patch}), segmenting based on whitespace~(\S\ref{section:white-patch}), and adaptive patching guided by entropy estimates from a compact byte-level language model~(\S\ref{section:dyn-patch}). We conclude by examining incremental patching and contrasting patching procedures with traditional tokenization methods~(\S\ref{section:bpe-patch}).

\subsection{Strided Patching Every K Bytes}
\label{section:static-patch}

A common method for aggregating bytes involves dividing them into uniformly sized patches of length $k$, as exemplified by MegaByte~\citep{yu2023megabyte}. This approach, characterized by its fixed stride, offers simplicity in both training and inference stages, allows for straightforward adjustment of the mean patch length, and thus facilitates easy regulation of computational expense in terms of flops. Despite these advantages, this strategy presents notable limitations. Primarily, computational resources are not adaptively assigned based on informational content: for instance, one may squander a transformer step $j$ on uninformative regions such as code whitespace, while failing to devote sufficient computation to areas rich in semantics, such as mathematical expressions. Additionally, such a scheme can result in inconsistent and context-insensitive partitioning of comparable byte sequences; for example, identical words may be fragmented differently depending on their positions.

\subsection{Space Patching}
\label{section:white-patch}

\citet{slagle2024spacebyte} introduces a straightforward but powerful enhancement to traditional strided patching by generating new patches whenever space-like bytes are encountered\footnote{
    Space-like bytes refer to bytes that are neither Latin characters, digits, nor UTF-8 continuation bytes; additionally, each patch must contain at least one byte that is not space-like.}. These bytes often coincide with linguistically meaningful segmentations in a wide range of languages. The Space patching approach devotes an extra latent transformer computation (i.e., increased flops) to each word, ensuring consistent patching for words across various sequences and prioritizing computational effort for segments—especially post-space segments—where predictive uncertainty is greatest. For instance, when answering ``Who composed the Magic Flute?\uline{\hspace{.6em}}'', predicting the initial byte of the answer is substantially more challenging than predicting subsequent bytes after ``M,'' because this first character dramatically narrows the set of possible completions, making the prediction of ``Mozart'' less uncertain for the remaining bytes. Nonetheless, Space patching does not generalize seamlessly to all languages or applications and inherently lacks the flexibility to adjust patch sizes. In the subsequent section, we present a novel patching technique motivated by the observation that initial bytes in words are typically hardest to forecast, while also introducing a flexible approach for patch size adjustment.

\subsection{Entropy Patching: Using Next-Byte Entropies from a Small Byte LM}
\label{section:dyn-patch}

Instead of depending on a rule-based heuristic like whitespace, we adopt a data-driven strategy to detect regions with elevated uncertainty in next-byte predictions. We propose \textit{entropy patching}, a method that utilizes entropy estimates to establish the limits of patches.

We train a small byte-level auto-regressive language model on the training data for {\textsc BLT} and compute next byte entropies under the LM distribution $p_{e}$ over the byte vocabulary $\mathcal{V}$:

\begin{equation}
H(x_i) = \sum_{v \in \mathcal{V}} p_{e}(x_i=v|\pmb{x}_{<i}) \log p_{e}(x_i=v|\pmb{x}_{<i})
\end{equation}

We investigate two approaches for detecting patch boundaries based on entropies $H(x_i)$. The first approach selects points exceeding a fixed, global entropy threshold, as shown in~\autoref{fig:patching}. The second method flags points whose entropy is substantially higher compared to the preceding value. Alternatively, this second strategy may be viewed as pinpointing locations where the trend of roughly monotonically decreasing entropy within the patch is disrupted.

\begin{align*}
    \text{Global Constraint}\hspace{1cm}H(x_t)& > \theta_g\\
    \text{Approx. Monotonic Constraint}\hspace{1cm}H(x_t)& - H(x_{t-1}) > \theta_r
\end{align*}

Patch boundaries are determined in a simple preprocessing phase carried out as part of the data loading process. This approach contrasts with that of~\citet{nawrot-etal-2023-efficient}, who employ a classifier specifically trained to predict patch boundaries based on entropy measurements. Within our experiments~(\S\ref{section:experiments}), we evaluate and contrast these two techniques for separating bytes of low and high entropy.

\subsection{The Byte-Pair Encoding (BPE) Tokenizer and Incremental Patching}
\label{section:incremental-patching}
\label{section:bpe-patch}

A large number of contemporary llms, such as our baseline Llama 3, employ subword tokenizers like bpe~\citep{gage1994new, sennrich-etal-2016-neural}. Here, we define ``tokens'' as byte sequences selected from a \textit{finite} vocabulary that is established before training begins, in contrast to ``patches,'' which denote dynamically constructed groups of bytes that are not constrained to a pre-defined vocabulary. A key distinction lies in the fact that, when utilizing tokens, the model lacks explicit access to the original byte-level information.

A significant advancement offered by {\textsc BLT} relative to tokenization-based models lies in its reconfiguration of the relationship between vocabulary size and computational requirements. In conventional large language models, augmenting the vocabulary size results in tokens that encode more bytes on average, thereby reducing the number of inference steps. However, this comes at the expense of enlarging the output dimension of the final projection layer, which increases memory and compute demands. As a result, tokenization-based methods are constrained in their capacity to meaningfully alter token granularity without incurring substantial changes in inference cost. For instance, Llama 3 achieves an increase in average token size from 3.7 to 4.4 bytes, but does so by quadrupling its embedding table size relative to Llama 2.

During generation, {\textsc BLT} must determine at each position in the byte sequence whether it lies at a patch boundary, as this choice dictates whether additional computation through the Latent Transformer will be triggered. Crucially, this decision must be made without relying on bytes that have not yet been generated. As a result, patching methods must not depend on knowledge of subsequent bytes to divide the byte sequence into patches. More precisely, a patching function $f_p$ is said to be incremental if it fulfills the following condition:
$$f_p(\pmb{x}_{<i}) = f_p(\pmb{x})_{<i}$$
It should be noted that bpe is not an incremental patching method, since the tokenization of a prefix may change depending on what follows, violating the stated property\footnote{Introducing an explicit delimiter token to mark patch boundaries can make bpe incremental, but this comes at the cost of longer byte sequences.}.

\section{{\textsc BLT} Architecture}
\label{section:architecture}
{\textsc BLT} comprises a sizeable global autoregressive language model functioning over patch representations. Additionally, it incorporates a pair of more lightweight local models: one for transforming byte sequences into patches and another for reconstructing byte sequences from patch representations~(\autoref{fig:arch_high_level}).

\subsection{Latent Global Transformer Model}

\textit{The Latent Global Transformer} refers to an autoregressive transformer architecture $\mathcal{G}$ composed of $l_{\mathcal{G}}$ layers, which transforms an input sequence of latent patch representations, $p_j$, into a corresponding sequence of output patch representations, $o_j$. In this work, we designate the subscript $j$ for patches, while $i$ corresponds to bytes. To ensure that attention is only paid up to the current patch within the current document, the global model incorporates a block-causal attention mask~\citep{dubey2024llama}. As this model accounts for the majority of computational operations, both during pre-training and at inference time, strategic decisions regarding its invocation enable us to modulate the computational resources required for different sections of the input and output, depending on their respective complexities.

\subsection{Local Encoder}
\label{section:local}

\textit{The Local Encoder Model}, represented by $\mathcal{E}$, is designed as a compact transformer leveraging significantly fewer layers ($l_{\mathcal{E}} << l_{\mathcal{G}}$) compared to its global counterpart. Its primary purpose is to transform an input sequence of bytes $b_i$ into highly descriptive patch representations $p_j$ with computational efficiency. Distinct from a standard transformer, this architecture integrates a cross-attention layer after each transformer block, enabling the aggregation of byte-level features into higher-level patch representations, as illustrated in~(\autoref{fig:crossattn}). 

Initially, each byte $b_i$ from the input is converted into a vector using an embedding matrix of dimension $\mathbb{R}^{256 \times h_{\mathcal{E}}}$, resulting in $x_i$. Optionally, these embeddings can be enriched by incorporating hash-embeddings~(\S\ref{sec:hash-emb}) that inject additional contextual cues. 

Next, a sequence of transformer and cross-attention layers~(\S\ref{sec:enc-cross-attn}) iteratively refines these byte-level representations, ultimately yielding patch representations $p_i$ suitable for downstream processing by the global transformer model $\mathcal{G}$. For attention, the transformer layers implement a \textit{local block causal} masking strategy: each byte may attend to a fixed span of $w_{\mathcal{E}}$ previous bytes, allowing for overlap across dynamic patch borders while strictly preventing attention beyond document edges. The ensuing subsections provide an in-depth discussion of the embeddings and the mechanics of the cross-attention module.

\subsubsection{Encoder Hash n-gram Embeddings}
\label{sec:ngram-emb}
\label{sec:hash-emb}
An essential aspect of forming resilient and informative representations at every step $i$ involves integrating context from prior bytes. In {\textsc BLT}, our approach captures this context by encoding both the single byte $b_i$ on its own and as a constituent within a byte n-gram structure. At each index $i$, we begin by assembling byte-grams

\begin{align}
    g_{i,n}=\{b_{i-n + 1},\ldots, b_i\}
\end{align}

for every byte position $i$ and for $n$ ranging from three to eight.\footnote{
Byte-grams of length $n$ or greater are excluded when $i<n$. }

Next, we present hash $n$-gram embeddings, which assign each possible byte $n$-gram, through a hash function, to a position within a fixed-size embedding matrix $E_{n}^{hash}$, with a separate matrix for each $n$ in $\{3,4,5,6,7,8\}$~\citep{bai2010learning}.
We then sum the embedding from $E_{n}^{hash}$ corresponding to the hashed $n$-gram with the embedding for the current byte, and this sum is normalized before it is input to the local encoder architecture. The computed embedding is as follows:

\begin{align}
e_i &= x_i + \sum_{n = 3,...,8}E_{n}^{hash}(\text{Hash}(g_{i,n})) \\
\text{where, Hash}(g_{i,n}) &= \text{RollPolyHash}(g_{i,n}) \% |E_{n}^{hash}|
\end{align}

We adjust $e_i$ by dividing it by the number of distinct $n$-gram sizes increased by one, utilizing the RollPolyHash method as specified in Appendix~\ref{appendix:rollpolyhash}. Section~\ref{section:ablations} presents an analysis in which we investigate the impact of varying $n$-gram hash embedding configurations—modifying both the value of $n$ and the embedding table size—on flop-controlled scaling law behavior. Beyond hash-based $n$-gram embeddings, our investigations also covered frequency-based $n$-gram embeddings, with further details regarding this line of inquiry documented in Appendix~\ref{appendix:ngrams}.

\subsubsection{Encoder Multi-Headed Cross-Attention}
\label{sec:enc-cross-attn}
Our approach largely mirrors the input cross-attention mechanism described in the Perceiver architecture~\citep{jaegle2021perceiver}. However, a key distinction is that, in our method, the latent representations map to variable patch representations rather than being a fixed latent set~(\autoref{fig:crossattn}), and they exclusively focus on the byte sequences within each individual patch. The construction of the module begins by forming a query vector for every patch $p_j$. Each query vector is created through pooling the byte embeddings of patch $p_j$ and subsequently applying a linear projection defined by $\mathcal{E}_{C} \in \mathbb{R}^{h_{\mathcal{E}} \times (h_{\mathcal{E}} \times U_{\mathcal{E}})}$, where $U_{\mathcal{E}}$ indicates the number of attention heads in the encoder cross-attention. To formalize, consider $f_{\text{bytes}}(p_j)$ as representing the byte sequence for patch $p_j$, and thus we compute

\begin{align}
P_{0,j} &= \mathcal{E}_{C}(f_\text{bytes}((p_j)), f~ \text{is a pooling function} \\
P_l &= P_{l-1} + W_o\left(\text{softmax}\left(\frac{QK^T}{\sqrt{d_k}}\right)V\right) \\
\text{where } Q_j &= W_q(P_{l-1,j}), K_i = W_k(h_{l-1,i}), V_i = W_v(h_{l-1,i}) \\
h_l &= \text{Encoder-Transformer-Layer}_l(h_{l-1})
\end{align}

In this context, $P \in \mathbb{R}^{n_p \times h_{\mathcal{G}}}$ denotes the set of $n_p$ patch representations that are intended for the global model. Each patch representation is generated by aggregating the corresponding byte embeddings $e_i$ within patch $p_j$. The matrices $W_q$, $W_k$, $W_v$, and $W_o$ are the learned projections associated with queries, keys, values, and output, respectively. Here, the keys and values are projected from the byte representations $h_i$ received from the preceding layer (with $e_i$ serving this purpose in the initial layer). For the attention mechanism, we implement a specialized masking approach whereby each query $Q_j$ is restricted to attend solely to keys and values linked to the bytes within the same patch $j$. Given that multi-head attention is utilized for $Q$, $K$, and $V$, and that patch representations have a typically larger dimensionality ($h_{\mathcal{G}}$) than $h_{\mathcal{E}}$, we keep $P_l$ structured as multiple heads, each of dimensionality $h_{\mathcal{E}}$, during cross-attention computations. These multiple head outputs are then concatenated to yield vectors in $\mathbb{R}^{h_{\mathcal{G}}}$. The architecture also applies pre-LayerNorm to queries, keys, and values before attention is computed, and omits the use of positional encodings in the cross-attention operation. A residual connection is included around the entire cross-attention module.

\subsection{Local Decoder} 

Analogous to the local encoder, the local decoder $\mathcal{D}$ utilizes a compact transformer architecture with significantly fewer layers, $l_{\mathcal{D}} << l_{\mathcal{G}}$. Its role is to transform a sequence of global patch representations, $o_j$, back into sequences of raw bytes, $y_i$. The decoding process involves generating each byte in the sequence conditioned on previously generated bytes, leveraging the hidden states provided by the local encoder corresponding to the byte sequence. The decoder sequentially applies $l_{\mathcal{D}}$ layers, where each iteration alternates between cross-attention and transformer operations. In each layer, the cross-attention mechanism is executed prior to the transformer, converting patch-level representations into corresponding byte-level representations, after which the transformer's operations are performed directly on the byte sequence.

\subsubsection{Decoder Multi-headed Cross-Attention} 

In the decoder cross-attention, the roles of the queries and key/values are interchanged i.e. the byte-representations are now the queries, and the patch representations are now the key/values. The initial byte-representations for the cross-attention are initialized as the byte embeddings from the last encoder layer i.e. $h_{l_{\mathcal{E}}}$. The subsequent byte-representations for layer $l$, $d_{l,i}$ are computed as:

\begin{align}
D_0 &= h_{l_{\mathcal{E}}} \\
B_l &= D_{l-1} + W_o\left(\text{softmax}\left(\frac{QK^T}{\sqrt{d_k}}\right)V\right), \\
  &\text{where } Q_i = W_q(d_{l-1,i}), K_i = W_k(\mathcal{D}_C(o_j)), V_i = W_v(\mathcal{D}_C(o_j)) \\
  D_l &= \text{Decoder-Transformer-layer}_l(B_l)
\end{align}

Here, $W_k$ and $W_v$ denote the key and value projection matrices, respectively; they are applied through a linear transformation followed by the split operation $\mathcal{D}_C$ to the global model's patch outputs $o_j$. The query projection, $W_q$, acts on the byte-level representations $d_{l-1}$ derived from the preceding decoder transformer layer—or, for the initial layer, $h_{l_{\mathcal{E}}}$. The projection matrix $W_o$ is applied to the outputs, resulting in $B \in \mathbb{R}^{h_{\mathcal{D}} \times n_b}$, where $n_b$ specifies the number of generated output bytes. The subsequent representations $D_l$ are produced by applying a decoder transformer layer to the output of the cross-attention module, $B$. Following the local encoder cross-attention design, multiple attention heads are utilized, LayerNorm is applied prior to submodules, no positional encodings are used, and a residual connection wraps the cross-attention component.

\section{Experimental Setup}
\label{section:experiments}
We establish a rigorous experimental protocol to objectively assess {\textsc BLT} relative to tokenization-based models, ensuring that {\textsc BLT} is not inadvertently advantaged through potentially longer accessible sequence contexts.

\subsection{Pre-training Datasets} In this study, models of all considered sizes undergo pre-training using two distinct datasets: (1) The Llama 2 dataset~\citep{touvron2023llama}, consisting of 2 trillion tokens sourced from a wide range of publicly accessible resources, which have been subsequently processed to remove noise and enhance dataset quality; and (2) {\textsc BLT}-1T: a newly constructed dataset with 1 trillion tokens drawn from diverse public datasets and supplemented with portions of pre-training data publicly released by Datacomp-LM~\citep{li2024datacomp}. The former dataset serves for scaling law investigations to identify the optimal token count as outlined by~\cite{dubey2024llama}, thereby informing architectural decisions for {\textsc BLT}; the latter facilitates a full-scale pre-training experiment, enabling direct downstream evaluation alongside Llama 3. Importantly, both collections exclude any content sourced from Meta’s own products or services. Additionally, when establishing baselines with tokenizer-dependent models, we employ the Llama 3 tokenizer configured with a 128K token vocabulary, since it consistently yielded better performance than the Llama 2 tokenizer throughout our benchmarks.

\subsection{Entropy Model} For our experiments, we utilize a byte-level language model as the entropy model, which is trained on the identical dataset as the complete {\textsc BLT} system. By default, unless specified otherwise, this model is a transformer architecture comprising 100 million parameters, structured over 14 layers, with each layer having a hidden size of 512. The model applies sliding window attention spanning 512 bytes. All other hyperparameters are kept consistent with those used in both our local and global transformer setups. We have explored various configurations regarding model scale, size of the receptive field, and alternate architectures, as outlined in \autoref{section:ablations}. Notably, when the model’s receptive field is sufficiently limited, it becomes feasible to represent the trained entropy model using an optimized lookup table.

\subsection{Entropy Threshold and Equalizing Context Length} For models that implement entropy-based patching, we determine a patching threshold so as to attain a target mean \textit{patch size} over the pretraining data blend. In contrast to tokenization, {\textsc BLT} allows for the arbitrary selection of \textit{patch size}, which can notably affect the model's utilized context length. To ensure that larger patch sizes do not result in an unintended benefit from longer contexts, we regulate the expected total number of bytes per batch, keeping it constant. As a consequence, for models with increased patch sizes, we proportionally decrease the sequence length. On Llama 2 data, a context window of 8k bytes is applied, whereas for the {\textsc BLT}-1T corpus, the context is extended to an average of 16k bytes, with the batch size held constant at an average of 16M bytes.

Although the average batch size remains unchanged, the proportion of bytes to patches varies across dynamic patching techniques when assembling data batches. To enhance efficiency, our {\textsc BLT} training implementation groups patches into batches, thereby preventing padding operations in the computationally intensive latent transformer component. This approach guarantees a uniform count of patches per batch. Throughout training, we apply padding and, when necessary, truncate byte sequences to 12k bytes for the Llama 2 dataset and 24k bytes for the {\textsc BLT}-1T dataset. This strategy mitigates memory usage surges that could arise from sequences containing particularly large patches.

\subsection{Entropy Model Context}
\label{section:entropy-context}
Our experiments indicate that applying entropy patching tends to produce increasingly larger patches, particularly in structured tasks such as multiple choice questions, where content repetition is common (refer to \autoref{fig:mmlu_patching} for an example on MMLU). These differences arise due to the reduced entropy associated with the repetitive content that the entropy model context captures. Therefore, in the large-scale deployment of {\textsc BLT}-Entropy with a patch size of 4.5, we refresh the entropy context using new line breaks and impose an approximate monotonicity constraint; this approach better mitigates "entropy drift" resulting from variation in context length. Notably, this adjustment solely impacts the calculation of entropies. The process for determining the entropy threshold value remains unchanged.

\subsection{FLOPs Estimation} 
\label{section:flops}

Our calculations for transformer floating point operations per second (flops) are primarily based on the methodology from Chinchilla~\citep{hoffmann2022training}, encompassing the feed-forward network, qkvo projections within the self-attention module, and the flops incurred during both attention calculation and output projection. Unlike Chinchilla, however, we presume that the input embedding layer is realized as an efficient lookup table rather than via dense matrix multiplication, resulting in this layer contributing zero flops. Consistent with established work, we approximate that the backward pass requires double the flops of the forward computation.

To compute flops \textit{per byte} for {\textsc BLT} models, we add up the flops for the local encoder transformer, the global latent transformer, and the local decoder transformer, together with the cross attention blocks in the encoder and the decoder:

\begin{align}
    \text{FL}_{\text{{\textsc BLT}}} &= \text{Transf. FL}(h_{\mathcal{G}}, l_{\mathcal{G}}, m=n_{ctx}/n_p, V=0) / n_p\\
   &\quad +\text{Transf. FL}(h_{\mathcal{E}}, l_{\mathcal{E}}, m=w_{\mathcal{E}}, V=0) \\
   &\quad + \text{Transf. FL}(h_{\mathcal{D}}, l_{\mathcal{D}}, m=w_{\mathcal{D}}, V=256) \\ 
    &\quad + \text{Cross Attn. FL}(h_{\mathcal{E}}, l_{\mathcal{E}}, m=n_p, r=n_p/k) \cdot k/n_p \\ 
   &\quad + \text{Cross Attn. FL}(h_{\mathcal{D}}, l_{\mathcal{D}}, m=k, r=k/n_p) 
\end{align}

Here, $n_{ctx}$ denotes the sequence length measured in bytes, $n_p$ specifies the patch size, $r$ represents the proportion of queries to key/value pairs, and $k$ defines the ratio of the patch dimension to the byte dimension—equivalently, it represents the number of separate local model partitions that are merged to create the overall model representation (see $k=2$ in \autoref{fig:crossattn}). The symbol $V$ indicates the output projection vocabulary size, which is relevant only in the context of the local decoder. The context length parameter, $m$, used by the attention mechanism varies based on whether it processes sequences of bytes or patches. We adjust the attention-related floating-point operation (flop) counts for each module accordingly. The precise formulas used for calculating Transformer-FLOPs and Cross-Attention FLOPs can be found in Appendix~\ref{section:flops-appendix}.

\subsection{Bits-Per-Byte Estimation} Perplexity only makes sense in the context of a fixed tokenizer as it is a measure of the uncertainty for each token. When comparing byte and token-level models, following previous work~\citep{xue2022byt5, yu2023megabyte, wang2024mambabyte}, we instead report Bits-Per-Byte (BPB), a tokenizer independent version of perplexity. Specifically:

Here, the cross-entropy loss summed over the data $\pmb{x}$ represents the model’s uncertainty, which is then scaled by dividing by both the number of bytes in $\pmb{x}$ and a logarithmic factor.

\subsection{Transformer Architecture Hyperparameters} In all transformer modules within {\textsc BLT}, encompassing both the local and global variants, we generally adhere to the configuration outlined for Llama~3~\citep{dubey2024llama}. Specifically, the feed-forward components employ the SwiGLU activation~\citep{swiglu}, while rotary positional embeddings (RoPE)~\citep{RoPe} with $\theta=500000$~\citep{xiong2024effective} are utilized exclusively in the self-attention layers. Layer normalization is performed using RMSNorm~\citep{RMSNorm}. For self-attention operations that implement fixed-pattern attention mechanisms—namely, \textit{block causal} and \textit{fixed-window block causal} masks—we apply Flash attention~\citep{dao2022flashattention}, setting the window size at 512 for fixed-width masking. In contrast, cross-attention layers require adaptive, patch-dependent attention masking; for these cases, we leverage Flex Attention\footnote{\url{https://pytorch.org/blog/flexattention}}, enabling efficient fused implementations that markedly accelerate model training.

\subsection{{\textsc BLT}-Specific Hyperparameters} 

To investigate the performance of {\textsc BLT} models, we perform two sets of experiments: analyzing scaling behaviors and assessing outcomes on downstream tasks. Our study covers models with varying sizes: 400M, 1B, 2B, 4B, and 8B, and corresponding architecture hyperparameters are detailed in Appendix~Table~\ref{tab:arch}. 

For initializing queries in the first cross-attention layer of the local encoder, we utilize max-pooling. All {\textsc BLT} models employ $500,000$ hashes using a single hash function, with n-gram lengths between 3 and 8. Across all scales, the learning rate is fixed at $4e-4$. This uniform learning rate selection for token-based and {\textsc BLT} models is informed by hyperparameter tuning at the 400M and 1B levels, where the search within the $1e-3$ to $1e-4$ range indicated that the same value was optimal. For experiments involving Llama-2 data, we adhere to batch-size recommendations from~\cite{dubey2024llama} or select equivalent batch sizes in terms of bytes. Optimization employs the AdamW algorithm~\citep{adamw}, setting $\beta_1$ to 0.9 and $\beta_2$ to 0.95, with $\epsilon=10^{-8}$. Learning rate scheduling includes a linear warm-up phase over 2000 steps, followed by cosine decay to zero. We implement weight decay at a rate of 0.1, and apply global gradient clipping with a cap of 1.0.

\section{Scaling Trends}
\label{section:scaling}

We offer a comprehensive analysis of the scaling patterns characteristic of byte-level models, providing guidance for further enlarging {\textsc BLT} models. This investigation is designed to overcome certain constraints present in earlier studies on byte-level models by: (a) examining the trends associated with compute-optimal training strategies, (b) training 8B-parameter models using substantial datasets (reaching as high as 1T tokens or 4T bytes) and assessing their effectiveness on downstream benchmarks, and (c) analyzing scaling behavior under conditions with fixed inference costs. Subsequent sections will explore unique benefits derived from processing byte-sequences.

\subsection{Parameter Matched Compute Optimal Scaling Trends}
\label{section:parameter-matched-scaling}
A range of \textit{compute-optimal} bpe and {\textsc BLT} models, spanning four parameter sizes from 1B to 8B, are trained on the Llama 2 dataset. Subsequently, we chart the language modeling accuracy against training flops for a representative portion of the training mixture. The bpe models utilize the parameter-to-data ratio considered optimal as established in Llama~3~\citep{dubey2024llama}. Such a \textit{compute-optimal} configuration is postulated, according to~\citep{hoffmann2022training}, to maximize model performance given a fixed computation budget, thus offering a strong comparative baseline. Additionally, for each bpe model, we construct a counterpart {\textsc BLT} model, trained on identical data and employing a Latent Transformer whose architecture and parameter count mirror those of the corresponding bpe Transformer.

As shown in~\autoref{fig:scaling-compute-opt} (right), {\textsc BLT} models consistently perform as well as, or better than, their bpe-based equivalents, and this pattern persists with increasing model size and computational budget. To our knowledge, {\textsc BLT} represents the first instance of a byte-level Transformer architecture exhibiting scaling behavior on par with BPE-based models in compute-optimal conditions. These findings support our hypothesis that the ideal parameter-to-compute ratio established for bpe models is also applicable to {\textsc BLT}, or at minimum, not substantially divergent.

Achieving alignment with bpe scaling behavior necessitates both advances in architecture and the use of dynamic patching. In ~\autoref{fig:scaling-compute-opt} (left), we present a comparison between models utilizing space-patching and Llama 3. Our approximation of SpaceByte \citep{slagle2024spacebyte} employs {\textsc BLT} space-patching, excluding n-gram embeddings and cross-attention. While SpaceByte yields improvements over Megabyte, its performance lags significantly behind that of Llama 3. In \autoref{fig:scaling-compute-opt} (right), we demonstrate the gains attributable to the combination of novel architectural choices and dynamic patching techniques. At large scale, {\textsc BLT} models achieve comparable performance to leading tokenizer-based approaches, including Llama 3.

Additionally, our results indicate that for models utilizing tokenizers, the selection of the tokenizer impacts performance: models trained with the Llama-3 tokenizer consistently achieve better results compared to those trained with the Llama-2 tokenizer, even when both are trained on identical datasets.

In summary, the {\textsc BLT} model exhibits behavior that falls between Llama 2 and Llama 3, especially when employing much larger patch sizes. While the average token size for the bpe tokenizers in Llama 2 and Llama 3 is 3.7 and 4.4 bytes, respectively, {\textsc BLT} maintains comparable scaling patterns when using average patch sizes of 6 or even 8 bytes. Because inference flop requirements are inversely related to average patch size, employing an 8-byte patch size could yield close to 50\% reduction in inference flop. Furthermore, as model and dataset sizes increase, models with these larger patch sizes tend to show improved performance. For example, although {\textsc BLT} with an 8-byte patch size initially underperforms relative to bpe Llama 2 at the 1B scale, it ultimately surpasses the bpe model by the 7B scale. This indicates that even greater patch sizes may deliver further benefits at even larger model and data scales as computational resources expand.

\subsection{Beyond Compute Optimal Task Evaluations}
\label{section:evals}
To further investigate scaling characteristics, we train an 8B {\textsc BLT} model at a scale that exceeds the compute-optimal regime, utilizing the {\textsc BLT}-1T corpus, which is both larger and of higher quality. Model performance is then assessed across a range of established classification and generation benchmarks. For evaluation, we cover tasks centered on common sense reasoning, general knowledge, and code generation:
\paragraph{Classification tasks} consist of ARC-Easy (0-shot)~\citep{Eval_ARC}, Arc-Challenge (0-shot)~\citep{Eval_ARC}, HellaSwag (0-shot)~\citep{Eval_hellaswag}, PIQA (0-shot)~\citep{Eval_piqa}, and MMLU (5-shot)~\citep{hendrycks2020measuring}. We adopt a prompt-based scoring approach, determining the probability distribution across the possible answer choices and presenting mean accuracy as the metric.
\paragraph{Coding related generation tasks:} We evaluate performance by reporting pass@1 on MBPP (3-shot)~\citep{Eval_MBPP} and HumanEval (0-shot)~\citep{Eval_HumanEval}, thereby measuring the capability of large language models to generate functional Python programs.

In \autoref{tab:evals}, we present a comparative analysis of three models, all trained on the {\textsc BLT}-1T dataset: a model utilizing the Llama 3 tokenizer with byte pair encoding (bpe),\footnote{We selected the Llama 3 tokenizer, featuring a 128k vocabulary, because it yields better results than the 32k vocabulary used in Llama 2.} as well as two {\textsc BLT} model configurations. The first variant incorporates a space-patching mechanism ({\textsc BLT}-Space), while the second adopts an entropy-guided patching approach ({\textsc BLT}-Entropy), where an approximate monotonicity constraint is applied, and the entropy model’s context is reset upon encountering new lines (as elaborated in~\autoref{section:entropy-context}). Each model was trained under comparable computational budgets (flop budget). Notably, with {\textsc BLT}-Entropy, we implement a modification during inference, lowering the entropy threshold from 0.6 to 0.1, which enhances task performance, albeit with an increase in the number of inference steps required.

The {\textsc BLT}-Entropy model achieves superior performance compared to the Llama 3 model on 4 of 7 evaluated tasks, despite being trained with an equal byte count. This enhanced performance likely arises from two factors: (1) more effective utilization of computational resources during training through dynamic patching, and (2) modeling information directly at the byte level rather than through tokenization.

Conversely, {\textsc BLT}-Space lags behind the Llama 3 tokenizer in performance across all but a single task. Nevertheless, due to its increased average patch size of 6 bytes, it notably reduces inference flops. For reference, both the bpe and entropy-patching models exhibit a similar mean patch size of about 4.5 bytes on the training dataset mix. With an equivalent training budget, the model with the larger patch size is able to process 30\% more data than the other two approaches, which could potentially move {\textsc BLT} further from the compute-optimal regime.

\subsection{Patches Scale Better Than Tokens} The {\textsc BLT} models allow for simultaneous scaling of both model capacity and patch dimension without surpassing the original computational constraints for training and inference, and while holding the training dataset unchanged. Unlike conventional token-based approaches bound by fixed vocabularies, patch-based models enable unbounded increases in patch size—granting a notable advantage in computational efficiency as detailed in Section~\ref{section:bpe-patch}. Employing larger patches reduces the required computations, permitting these saved resources to be invested in enlarging the global latent transformer, since its invocation frequency decreases as patch size grows.

We perform a fixed inference scaling analysis to evaluate whether increasing model size while reducing the number of steps on larger patches outperforms using smaller models with more steps. Beginning with models of 400m and 3.6B parameters and employing the Llama 2 tokenizer, we identify flop-matched counterparts utilizing the Llama 3 tokenizer, as well as {\textsc BLT}-Entropy models configured with mean patch sizes of 6 and 8 bytes on the training datamix (refer to~\autoref{tab:fixed-inf-params} for specifics regarding the models). For models with a patch size of 8, three encoder layers are used rather than one. Each model is trained with a range of training flop allocations.

As illustrated in \autoref{fig:fixed_inference_scaling}, {\textsc BLT} models exhibit superior scaling behavior compared to tokenization-based models across both inference flop categories. For both settings, BPE models initially outperform when training resources are limited, yet are soon overtaken by {\textsc BLT} models once the training allocation moves slightly beyond the compute-optimal threshold. In real-world scenarios, it is often advantageous to invest more heavily in the initial pretraining phase in order to produce a model that yields improved outcomes for a given inference budget. An apt illustration of this phenomenon is provided by the 8B model class—such as Llama 3.1—which was pretrained on data volumes approximately two orders of magnitude greater than the compute-optimal scale for its model size.

The point at which {\textsc BLT} surpasses token-based architectures has moved closer to the compute-optimal threshold with the increase to larger flop class models—shifting from approximately 3 times down to 2.5 times the optimal compute allocation. Additionally, within the higher flop class, the model utilizing a patch size of 8 exhibits a sharper scaling trajectory, surpassing other configurations at an earlier stage. As elaborated in Section~\ref{section:parameter-matched-scaling}, larger patch sizes demonstrate performance approaching that of BPE models as model size grows. We partially attribute this to the reduced proportion of total flops allocated to the byte-level Encoder and Decoder components, whose scaling lags behind that of the Latent Transformer. Notably, as the parameter count is increased 20-fold from 400M to 8B, the local parameter count in {\textsc BLT} increases by only about a factor of two. This distinction is critical, since only the patch Latent Transformer’s computational cost is influenced by larger patch sizes, leaving the flops used by byte-level modules largely unaffected. This observation explains why the relative size of {\textsc BLT}-Entropy with ps=8 increased from 1.6x to 1.7x of the Llama 2 parameter count when scaling up the model.

In conclusion, our investigation into patch-length scaling reveals that the {\textsc BLT} architecture benefits from jointly enlarging both patch size and model size, resulting in more favorable scaling behavior. Notably, these positive scaling dynamics tend to hold and can even be enhanced as the model size increases further.

\section{Byte Modeling Improves Robustness} We further evaluate the robustness of {\textsc BLT} against token-based models that do not incorporate byte-level representations, and introduce a technique for transforming existing pretrained token-based models to process byte information directly.
\label{sec:robustness}

\subsection{Character-Level Tasks}

One of the initial reasons for developing byte-level models was their resilience to noise introduced at the byte level in input data and their capacity to recognize the compositional structure within tokens—an area where conventional tokenizer-based approaches often fall short. To investigate these properties, we conduct supplementary tests utilizing benchmarks specifically designed to assess resistance to input noise as well as sensitivity to individual characters. These tests cover both English and multiple other languages, encompassing not only characters but also digits and phonemes. The outcomes of these evaluations are summarized in Table~\ref{tab:char_tasks}.

\paragraph{Noisy Data} To assess the comparative robustness of tokenizer-based approaches and {\textsc BLT}, we generate noisy variants of the benchmark classification datasets referenced in Section~\ref{section:evals}. Five unique character-level noise methods are utilized to introduce distortions into the input: (a) \textit{AntSpeak}: All characters are converted to uppercase and separated by spaces; (b) \textit{Drop}: Random deletion of 10\% of characters from the input; (c) \textit{RandomCase}: Half of all characters are randomly assigned to either uppercase or lowercase; (d) \textit{Repeat}: 20\% of the characters are each repeated up to four times; (e) \textit{UpperCase}: The text is entirely uppercased. For evaluation, each noise procedure is individually applied to the prompt, the completion, or both, constituting separate test cases, and the final scores are averaged. Table \ref{tab:char_tasks} presents performance on noisy versions of the HellaSwag dataset~\citep{Eval_hellaswag}, demonstrating that {\textsc BLT} consistently surpasses tokenizer-based models in robustness. On average, it achieves an 8-point lead over its tokenizer-trained counterpart and even yields better results than the Llama 3.1 model trained on a substantially larger corpus.

\paragraph{Phonology - Grapheme-to-Phoneme (G2P)} We evaluate {\textsc BLT}'s proficiency in converting sequences of graphemes (characters that represent words) into their corresponding phonemic transcriptions (the sounds that make up the word’s pronunciation). Table \ref{tab:char_tasks} displays performance metrics for the G2P task in a 5-shot framework, utilizing the Phonology Bench~\citep{suvarna2024phonologybench}. The results indicate that {\textsc BLT} achieves superior performance compared to the Llama 3 1T tokenizer-based model baseline on this task.

\paragraph{CUTE}
To evaluate character-level comprehension, we test {\textsc BLT} using the CUTE benchmark~\citep{edman2024cute}, which features multiple tasks grouped into three main areas: compositional understanding, recognition of orthographic similarity, and sequence manipulation skills. This suite of tasks proves especially difficult for most models reliant on tokenizers, as these systems seem aware of their token spellings but encounter challenges when attempting to leverage this information to manipulate text effectively. As reported in Table~\ref{tab:char_tasks}, {\textsc BLT}-Entropy surpasses both BPE Llama 3 variants by over 25 points on the benchmark. Notably, our model achieves near-perfect results in character manipulation, obtaining a 99.9\% success rate in both spelling-related evaluations. These substantial gains, even though {\textsc BLT} is trained with 16 times less data compared to Llama 3.1, suggest that acquiring character-level capabilities is inherently more difficult for BPE models. Figure \ref{fig:cute_examples} presents example cases where the Llama 3 tokenizer model underperforms, yet {\textsc BLT} delivers strong results. Word-level deletion and insertion tasks are the sole instances in which BPE outperforms; byte-level approaches find such manipulations less straightforward, though the margin is slim, and it may be easier for character-based architectures to learn word-level operations than for token-based approaches to acquire detailed character-level knowledge. Across all tasks, we maintain a consistent evaluation protocol, relying on Huggingface's original prompts. Enhanced prompt engineering could potentially further aid BPE models in these settings.

\paragraph{Low Resource Machine Translation} We assess the performance of {\textsc BLT} on both translation to and from English across six major language families and twenty-one low-resource languages—encompassing diverse writing systems—utilizing the FLORES-101 benchmark~\citep{goyal2022flores}. SentencePiece BLEU scores, as detailed in Table~\ref{tab:flores}, summarize these findings. The evaluation reveals that {\textsc BLT} surpasses a system built with the Llama 3 tokenizer, with an observed improvement of 2 BLEU points when translating into English, and a 0.5-point gain when translating from English. For widely studied language pairs, {\textsc BLT} matches or modestly exceeds the results of Llama 3. Notably, in the context of many low-resource language families, {\textsc BLT} achieves clear gains over Llama 3, highlighting the advantages of byte-level modeling in handling infrequent or long-tail byte sequences.

\subsection{Training {\textsc BLT} Using Llama 3}
\label{sec:distillation}
We investigate a method in which {\textsc BLT} models benefit from the initialization with parameters from existing pretrained tokenizer-based models, resulting in more efficient and accelerated convergence during training. Specifically, we initialize the global transformer parameters of {\textsc BLT} with those from a pretrained Llama 3.1 model. During the subsequent training, the global transformer’s weights are updated at a learning rate one-tenth of that used for the local encoder and decoder, following the recommended step schedule for Llama 3. We report results in Table~\ref{tab:distill}, comparing this approach to a baseline {\textsc BLT}. The findings show that {\textsc BLT} initialized from Llama 3.1 achieves markedly better performance than both the original Llama 3 and baseline {\textsc BLT}, while consuming the same number of compute flops. Additionally, despite the {\textsc BLT}-Entropy model (referenced in Table~\ref{tab:evals}) being trained on a much larger dataset of 1T tokens, {\textsc BLT} from Llama 3.1 surpasses it on the MMLU task. These results indicate that this strategy may substantially reduce training compute while delivering strong performance.

This approach can be interpreted as converting models that rely on tokenizers into ones that do not, thereby transforming a pre-trained LLaMA 3.1 model into a {\textsc BLT} framework. For a thorough comparison, Table \ref{tab:distill} reports results for the original LLaMA 3.1, which was pre-trained on 15T tokens, alongside the {\textsc BLT} version constructed from LLaMA 3. When evaluating these models, ours shows a modest reduction in performance on MMLU and HumanEval, while losses are more pronounced across other benchmark tasks. These outcomes indicate that additional research is necessary to more effectively utilize the pre-trained model, with an emphasis on refining data mixtures and optimizing hyperparameters to enhance performance.

\section{Ablations and Discussion}
\label{section:ablations}
Here, we present ablation studies that validate our design decisions regarding the architecture of {\textsc BLT}, as well as the selection of patching methodology and hyper-parameter settings used for the {\textsc BLT} model with 8B parameters, which was trained using the {\textsc BLT}-1T dataset.

\paragraph{Entropy Model Hyper-parameters} To investigate how entropy model capacity and context window size influence scaling behavior, we conduct experiments with byte-level entropy transformer architectures ranging from 1 million to 100 million parameters and context window lengths set between 64 and 512 tokens. We generate bpb versus training flop scaling law plots using our $400m$ and $1b$ {\textsc BLT} models—both trained on the Llama-2 dataset—and display these results in \autoref{fig:entropy_ablation}. The results indicate a positive association between scaling performance and increases in both entropy model size and context window length; however, returns start to plateau as model size exceeds 50 million parameters.

\paragraph{Types of Patching}

We conduct an ablation study on the four patching strategies detailed in Section~\ref{section:patching}: (1) Strided Patching with strides set to 4 and 6, (2) Patching at whitespace locations, (3) BPE Tokenizer patching utilizing the Llama 3 tokenizer, and (4) Entropy-driven patching employing a compact byte-level language model.

Although dynamic patching leads to shorter effective sequence lengths, we standardize sequence length across all patching methods to ensure that each one operates over comparable context lengths. During both training and inference, all models are exposed to an equivalent expected number of bytes per sequence, thereby eliminating confounding variables related to modeling longer contexts. As depicted in Figure~\ref{fig:scaling-compute-opt}, these ablation studies demonstrate that every alternative patching approach surpasses static patching, with space patching exhibiting performance nearly equivalent to that of dynamic entropy-based patching.

In \autoref{tab:evals_ablation}, we report benchmark results for {\textsc BLT} models, comparing tokenizer-based approaches, space patching, and entropy-based patching methods. All models are trained on the Llama 2 dataset for an optimal number of training steps~\citep{dubey2024llama}. While space patching employs a more straightforward approach that avoids invoking the entropy model during training, our analysis reveals that the improvements shown by entropy-based patching in scaling experiments~(Section~\ref{section:scaling}) are also reflected in performance on downstream benchmark tasks.\footnote{Space patching figures are from earlier runs without cross-attention, but the observed trends persist when cross-attention is included.}

\paragraph{Cross-Attention}
In~\autoref{tab:crossattn_ablations_1b}, we investigate the impact of inserting cross-attention modules at different locations within both the encoder and decoder components of {\textsc BLT} through an ablation study. Regarding encoder cross-attention, our experiments involve initializing the queries in three distinct manners: (1) utilizing a shared learned embedding for all global states, (2) employing a hash-based embedding derived from the patch’s byte content, and (3) applying a pooling operation to the encoder’s hidden states for the patch bytes at a specific encoder layer.

Our results indicate that implementing cross-attention within the \textit{decoder} yields the greatest benefit. In contrast, integrating cross-attention into the encoder offers a modest gain, observed primarily when queries are initialized through pooling. Furthermore, our analysis reveals that the advantages of cross-attention are more pronounced on Common-Crawl data, and this effect is amplified when working with larger patch sizes.

\paragraph{n-gram Hash Embeddings} 

We examine the impact of varying n-gram hash embedding vocabulary sizes—specifically 0, 100K, 200K, and 400K—and summarize the findings in Table~\ref{tab:ngram_ablations_1b}. Our results reveal that introducing hash embeddings yields benefits across all evaluated domains, with notable improvements on Wikipedia and Github (a 0.04 bpb gain compared to just 0.01 bpb after 15k steps at the 8B scale). Reducing the hash count from 500K to 300K at 8B parameters leads to only a marginal performance decrease of approximately 0.001 bpb over 15k steps. These observations suggest that hash embeddings are crucial for enabling {\textsc BLT} to reach parity with tokenizer-based architectures. However, increasing the hash vocabulary beyond 300K confers minimal additional advantage. Furthermore, the enhancements from hash embeddings and cross-attention tend to be complementary, as they improve outcomes on different types of datasets.

\paragraph{Local Model Hyperparameters} Table \ref{tab:local_layers} presents an ablation study examining different configurations for the number of layers within the local encoder and decoder. When utilizing hash n-gram embeddings, {\textsc BLT} demonstrates strong performance with a highly compact encoder architecture—specifically, one consisting of only a single layer—while benefitting from a more substantial decoder.

\section{Related Work}
\label{section:related_work}

\textbf{Character-Level RNNs:} Character Language Modeling has been widely explored since the inception of neural approaches~\citep{sutskever2011generating,mikolov2012subword,graves2013generating}, largely due to its inherent ability to represent out-of-vocabulary terms naturally, thus removing dependence on back-off techniques. \cite{kim2016character} introduce a model that solely handles character inputs via convolutional and highway layers, subsequently connecting to LSTM-based RNNs. Their approach achieves comparable results to then-leading RNN-based language models for English and exhibits superior performance for morphologically complex languages, highlighting a key strength of character-level LLMs. Byte-level LSTM models are applied to machine comprehension by \cite{kenter2018byte}, who report performance gains over word-level models for languages such as Turkish and Russian, both rich in morphology. Similarly, \cite{zhang2015character} employ convolutional neural networks over character sequences for classification purposes, achieving improved outcomes relative to word-level baselines on specific tasks. Hierarchical LSTM architectures featuring boundary detectors, as proposed by \citet{chung2019hierarchical}, learn latent text structure across multiple levels, thereby enhancing results in character-level language modeling. In another direction, ByteNet~\cite{kalchbrenner2016neural} adopts a convolutional approach at the character level for machine translation tasks, foregoing attention mechanisms in favor of CNN layers.

\textbf{Character-Level Transformers:} The introduction of transformer architectures leveraging attention mechanisms~\citep{vaswani2017attention} and subword tokenization schemes~\citep{sennrich-etal-2016-neural} yielded marked advances in language modeling as well as various benchmark evaluations. Yet, both word and subword representations inherently encode preferences regarding the appropriate abstraction for model reasoning. Seeking to integrate the strengths of transformers with early achievements in character-level modeling, \citet{al2019character} applied exceptionally deep transformer networks and, utilizing supplementary auxiliary losses, managed to surpass earlier LSTM-based character language models. Nevertheless, a notable disparity persisted relative to word-level LLMs. Similarly, GPT-2~\citep{radfordgpt2} found that, in large-scale corpora such as the 1 Billion Word Benchmark, models trained at the byte level underperformed compared to their word-level counterparts.

The work of \cite{Choe2019BridgingTG} illustrated that transformer-based LLMs operating at the byte level can surpass their subword-level counterparts in performance for models with similar parameter counts; however, these architectures incur substantially higher computational costs and prolonged training durations. Likewise, \cite{el2020characterbert} introduced CharFormer, a BERT variant that generates word representations via convolutions applied on character embeddings. This model delivered notable improvements within the medical domain, albeit at a significantly increased computational expense. The CANINE model introduced by \cite{clark2022canine}, comprised of 150M parameters and restricted to an encoder-only design, processes text directly as character sequences. Its architecture features a deep transformer backbone, analogous to the global model in this work, and employs a local transformer module together with strided convolutional layers for input character downsampling. CANINE demonstrated superior downstream multilingual task performance relative to token-level models such as mBERT. ByT5~\citep{xue2022byt5} advanced byte-level encoder-decoder frameworks that avoided patching techniques altogether; although their design achieved greater resistance to noisy data and matched or exceeded the performance of tokenized models with significantly less data (by a factor of four), the absence of patching operations resulted in the need for intensive attention computation at the byte level, thus dramatically increasing compute demands. Representing sequences at the byte level, as opposed to using subword units, not only prolongs input sequence lengths but also complicates efficient scaling of such models. More recently, employing the Mamba Architecture~\citep{gu2023mamba}—capable of sustaining fixed-size memory across very extended contexts—\cite{wang2024mambabyte} succeeded in training a byte-level Mamba model without resorting to patching operations. Their architecture achieved state-of-the-art results, outperforming byte-level transformers in flop-controlled settings at the 350M parameter scale, as evidenced by improved bits-per-byte scores on several benchmark datasets.

\textbf{Patching-based approaches:} Leveraging patching strategies has the potential to significantly reduce the excessive computation (in terms of flops) characteristic of byte-level LLMs, while possibly preserving or even enhancing their performance. Several studies have showcased promising results, particularly when applied to smaller model architectures and modest quantities of training bytes. For instance, \cite{nawrot-etal-2022-hierarchical} introduced the hourglass transformer by incorporating static patching through both downsampling and upsampling operations, demonstrating superior results compared to other byte-level baselines at the 150M parameter level. Expanding on this, \cite{nawrot-etal-2023-efficient} advanced the approach via dynamic patching strategies, which included an end-to-end trainable boundary predictor, a variant of the boundary predictor utilizing supervision from existing tokenizers, as well as an entropy-driven patching method analogous to {\textsc BLT}. Their findings indicated that this dynamic patching methodology can surpass the performance of conventional transformers of the period, particularly in language modeling, at a scale of 40M parameters and using approximately 400M tokens. Separately, \cite{lester2024training} explored compressing sequences with arithmetic coding, thus achieving compression levels exceeding those attainable by BPE. Applying an equal-info windows approach, their system surpassed byte-level baselines on language modeling benchmarks, although it did not outperform subword-based baselines.

Our research is heavily influenced by, and bears the closest resemblance to, MegaByte~\citep{yu2023megabyte}, a causal, decoder-only language model that employs a fixed, static patching scheme. In MegaByte, byte sequences are transformed into patches by concatenating their representations, and a localized decoder-side module reconstructs the original bytes from these patches. The original MegaByte work shows that this model can achieve parity with tokenizer-based language models containing 1B parameters, evaluated on a data corpus of roughly 400B bytes. We include comprehensive ablations of MegaByte in our studies, finding that the static patching approach underperforms state-of-the-art tokenizer-based models that are optimized for computational efficiency when training resources are tightly controlled; we further demonstrate how {\textsc BLT} effectively closes this performance gap. Similarly, \cite{slagle2024spacebyte} report analogous findings regarding MegaByte, and propose adaptations such as applying patching to whitespaces and similar byte classes, in addition to introducing a localized encoder module. Their analysis reveals performance gains over tokenized transformer models in compute-limited regimes for specialized datasets like Github and arXiv, specifically at the 1B parameter level. We extend our evaluation to include this architecture, concluding that additional innovations in model design are essential to successfully scale byte-level models to match, or surpass, the capabilities of leading token-based systems such as Llama 3.

\section{Limitations and Future Work}

In this study, our selection of architectural configurations is based on training models for the number of steps found to be optimal for Llama~3~\citep{dubey2024llama}. It is important to note, however, that these scaling laws were originally established for BPE-level transformer architectures and thus may not produce ideal ratios between data and parameter sizes when applied to {\textsc BLT}. Deriving precise scaling laws specifically tailored to {\textsc BLT}—which might yield even more advantageous scaling properties for our model—is an avenue we intend to explore in future research. Furthermore, the majority of our experimental analysis has been limited to models with up to 1B parameters. Architectural optimality could shift as we increase model sizes to 8B parameters or higher, potentially resulting in superior performance at these larger scales.

Current transformer libraries and codebases prioritize efficiency for architectures based on tokenization. Although our study offers theoretical flop-equivalent comparisons and incorporates some optimized solutions (for example, FlexAttention) for layers that differ from standard transformer structures, our practical implementations may still lag behind tokenizer-centric models with respect to wall-clock runtime, indicating potential gains from additional optimization efforts.

Whereas {\textsc BLT} employs an independently trained entropy model to perform patching, an avenue for future research involves integrating the patching model into a fully end-to-end learning paradigm. In Section \ref{sec:distillation}, we provide preliminary experimental results that suggest the viability of “byte-ifying” models such as Llama 3—which utilize tokenizers and have been trained on over 10T tokens—by leveraging initialization of the global transformer with these pre-existing weights and keeping them fixed during training. Advancing research along this pathway could yield techniques that preserve the advantages conferred by bytefying, and potentially enhance performance to levels exceeding those achieved by the original tokenizer-based models, all without necessitating training from the ground up.

\section{Conclusion}
\label{section:conclusion}

This work introduces the Byte Latent Transformer (\textbf{BLT}), an innovative model architecture that challenges the standard reliance on fixed-vocabulary tokenization in large language models. By employing a dynamic and trainable approach to aggregating bytes into variable-length patches, {\textsc BLT} enables more flexible allocation of computational effort in response to data complexity, which yields notable enhancements in efficiency and resilience. Through a comprehensive scaling analysis, we show that {\textsc BLT} models are able to achieve parity with tokenization-dependent systems such as Llama 3 for model sizes up to 8B parameters and datasets of 4T bytes, with the potential to attain up to a 50\% reduction in inference FLOPs at the cost of only minor decreases in evaluation metrics. Additionally, {\textsc BLT} offers a unique scaling pathway, permitting joint escalation of both model size and patch size while maintaining a constant inference cost—a beneficial property for the computation constraints prevalent in applied environments. By working directly on byte-level inputs, {\textsc BLT} not only strengthens robustness against irregular or noisy data but also enhances the representation of sub-word elements, affording improved understanding of rare patterns. Collectively, these advancements establish {\textsc BLT} as a compelling, scalable, and flexible alternative to conventional tokenization strategies for building more efficient and robust language models.

\end{document}